\usepackage[T1]{fontenc}
\usepackage{babel}
\usepackage{graphics}
\usepackage{geometry}
\usepackage{epsfig}
\usepackage{mathptmx}
\usepackage{amsmath}
\usepackage{amssymb}
\usepackage{bbm}
\usepackage{datetime}
\usepackage{color}
\usepackage{physics}
\usepackage{siunitx}
\usepackage{booktabs}
\usepackage{pgfplots}
\usepackage{tikz}
\usepackage{mathtools}
\usepackage{hyperref}
\usepackage{cleveref}
\usepackage{tikz}
\usepackage{pgfplots}
\usepackage{enumitem}
\usepackage{subcaption}
\usepackage{multirow}
\usepackage{minitoc}
\usepackage{titlesec}
\usepackage[style=authoryear-comp]{biblatex}

%%%%%%%%%%%%%%%%%%%%%%%%%%%%%%%%%%%%%%%%%%%%%%%%%%%%%%%%%%%%

\addbibresource{bib.bib}

%%%%%%%%%%%%%%%%%%%%%%%%%%%%%%%%%%%%%%%%%%%%%%%%%%%%%%%%%%%%

\AtBeginDocument{\RenewCommandCopy\qty\SI}

%%%%%%%%%%%%%%%%%%%%%%%%%%%%%%%%%%%%%%%%%%%%%%%%%%%%%%%%%%%%

\graphicspath{{figures/}}

\usetikzlibrary{external, angles, quotes, calc, positioning, pgfplots.statistics, patterns}

\newdimen\figurewidth
\figurewidth=0.7\textwidth

\tikzset{external/system call={
			pdflatex
			\tikzexternalcheckshellescape
			--extra-mem-top=10000000 % that helps compile humongus tikz pictures
			--main-memory=10000000
			--extra-mem-bot=10000000
			-interaction=batchmode
			-jobname
			"\image"
			"\texsource"
		}}

\tikzexternalize[prefix=build/]    % PAS DE '{}' À CETTE COMMANDE CA CASSE TOUT

\newcommand{\inputtikz}[1]{
	\tikzsetnextfilename{#1}
	\input{#1}
}

%%%%%%%%%%%%%%%%%%%%%%%%%%%%%%%%%%%%%%%%%%%%%%%%%%%%%%%%%%%%

\dominitoc
\tightmtctrue
\setcounter{minitocdepth}{1}

\titleformat{\chapter}[display]
{\normalfont\huge\bfseries}{\chaptertitlename\ \thechapter}{20pt}{\Huge\hyperlink{table_of_contents}}

\titleformat{\section}
{\normalfont\Large\bfseries}{\thesection}{1em}{\hyperlink{table_of_contents}}

\titleformat{\subsection}
{\normalfont\large\bfseries}{\thesubsection}{1em}{\hyperlink{table_of_contents}}

\titleformat{\subsubsection}
{\normalfont\normalsize\bfseries}{\thesubsubsection}{1em}{\hyperlink{table_of_contents}}

\newcommand\shorttableofcontents{
	{
			\renewcommand*{\contentsname}{Short contents}
			\setcounter{tocdepth}{0}
			\tableofcontents*
		}
}
\newcommand\longtableofcontents{
	{
			\renewcommand*{\contentsname}{Full Contents}
			\setcounter{tocdepth}{10}
			\hypertarget{table_of_contents}{}
			\tableofcontents
		}
}

%%%%%%%%%%%%%%%%%%%%%%%%%%%%%%%%%%%%%%%%%%%%%%%%%%%%%%%%%%%%

\nouppercaseheads

%%%%%%%%%%%%%%%%%%%%%%%%%%%%%%%%%%%%%%%%%%%%%%%%%%%%%%%%%%%%

% \crefformat{subsection}{Section~#2#1#3}
% \Crefformat{subsection}{Section~#2#1#3}

% \crefformat{equation}{(#2#1#3)}

% \crefformat{table}{Tab.~#2#1#3}
% \Crefformat{table}{Tab.~#2#1#3}

% \crefformat{figure}{Fig.~#2#1#3}
% \crefrangeformat{figure}{Fig.~#3#1#4 to~#5#2#6}
% \crefmultiformat{figure}{Fig.~#2#1#3}{ and~#2#1#3}{,~#2#1#3}{ and~#2#1#3}
% \Crefformat{figure}{Fig.~#2#1#3}
% \Crefrangeformat{figure}{Fig.~#3#1#4 to~#5#2#6}
% \Crefmultiformat{figure}{Fig.~#2#1#3}{ and~#2#1#3}{,~#2#1#3}{ and~#2#1#3}

%%%%%%%%%%%%%%%%%%%%%%%%%%%%%%%%%%%%%%%%%%%%%%%%%%%%%%%%%%%%

\DeclareMathOperator*{\argmin}{arg\!\min}

\newcommand\coord[2]{\ensuremath{{}^{#2}#1}}
\newcommand\lucoord[3]{\ensuremath{{}^{#2}#1^{#3}}}
\newcommand\ldcoord[3]{\ensuremath{{}^{#2}#1_{#3}}}

\newcommand\point[1]{\ensuremath{\mathbf{#1}}}
\newcommand\fpoint[2]{\ensuremath{{}^{#2}\point{#1}}}
\newcommand\lupoint[2]{\ensuremath{\point{#1}^{#2}}}
\newcommand\ldpoint[2]{\ensuremath{\point{#1}_{#2}}}
\newcommand\flupoint[3]{\ensuremath{\fpoint{#1}{#2}^{#3}}}
\newcommand\fldpoint[3]{\ensuremath{\fpoint{#1}{#2}_{#3}}}
\newcommand\fllpoint[4]{\ensuremath{\fpoint{#1}{#2}_{#3}^{#4}}}

\newcommand\vect[1]{\ensuremath{\vec{\mathbf{#1}}}}
\newcommand\fvect[2]{\ensuremath{{}^{#2}\vect{#1}}}
\newcommand\luvect[2]{\ensuremath{\vect{#1}^{#2}}}
\newcommand\ldvect[2]{\ensuremath{\vect{#1}_{#2}}}
\newcommand\fluvect[3]{\ensuremath{\fvect{#1}{#2}^{#3}}}
\newcommand\fldvect[3]{\ensuremath{\fvect{#1}{#2}_{#3}}}
\newcommand\fllvect[4]{\ensuremath{\fvect{#1}{#2}_{#3}^{#4}}}

\newcommand\res[1]{\ensuremath{\varepsilon_{#1}}}

\newcommand\lplane[1]{\ensuremath{\mathcal{P}_{#1}}}
\newcommand\lframe[1]{\ensuremath{\mathcal{F}_{#1}}}

\newcommand\fpointfull[2]{\ensuremath{\fpoint{#1}{#2}=\left(\coord{x}{#2}, \coord{y}{#2},\coord{z}{#2}\right)}}
\newcommand\flupointfull[3]{\ensuremath{\flupoint{#1}{#2}{#3}=\left(\lucoord{x}{#2}{#3}, \lucoord{y}{#2}{#3},\lucoord{z}{#2}{#3}\right)}}
\newcommand\fldpointfull[3]{\ensuremath{\fldpoint{#1}{#2}{#3}=\left(\ldcoord{x}{#2}{#3}, \ldcoord{y}{#2}{#3},\ldcoord{z}{#2}{#3}\right)}}

\newcommand\fvectfull[2]{\ensuremath{\fvect{#1}{#2}=\left(\coord{x}{#2}, \coord{y}{#2},\coord{z}{#2}\right)}}
\newcommand\fluvectfull[3]{\ensuremath{\fluvect{#1}{#2}{#3}=\left(\lucoord{x}{#2}{#3}, \lucoord{y}{#2}{#3},\lucoord{z}{#2}{#3}\right)}}
\newcommand\fldvectfull[3]{\ensuremath{\fldvect{#1}{#2}{#3}=\left(\ldcoord{x}{#2}{#3}, \ldcoord{y}{#2}{#3},\ldcoord{z}{#2}{#3}\right)}}

\newcommand\lplanefullxy[1]{\ensuremath{\lplane{#1}=\left(\ldpoint{O}{#1}, \ldvect{x}{#1},\ldvect{y}{#1}\right)}}
\newcommand\lplanefullxz[1]{\ensuremath{\lplane{#1}=\left(\ldpoint{O}{#1}, \ldvect{x}{#1},\ldvect{z}{#1}\right)}}
\newcommand\lplanefullyz[1]{\ensuremath{\lplane{#1}=\left(\ldpoint{O}{#1}, \ldvect{y}{#1},\ldvect{z}{#1}\right)}}

\newcommand\lframefull[1]{\ensuremath{\lframe{#1}=\left(\ldpoint{O}{#1},\ldvect{x}{#1},\ldvect{y}{#1},\ldvect{z}{#1}\right)}}

\newcommand\inzeroton{\ensuremath{\in\{0,\hdots, n\}}}

\newcommand\thetamodel{$\theta$-augmented}
\newcommand\gammamodel{$\gamma$-augmented}
\newcommand\thetagammamodel{$\theta\gamma$-augmented}